\chapter{攻击方法设计与实现}
\label{cha:method}

本章给出AES Flush+Reload攻击测试框架的完整设计与实现方案。攻击框架分两个阶段协同工作：第一阶段通过Flush+Reload侧信道监控T表缓存行的访问模式，利用排除法恢复最后一轮子密钥；第二阶段复用同一批监控数据，结合AES等效解密的代数推导恢复倒数第二轮子密钥，经密钥扩展逆运算反解出原始密钥。针对AES-128、AES-192和AES-256三种密钥长度的差异，框架采用统一的T表监控策略和差异化的密钥恢复流程：AES-128仅需第一阶段即可完成攻击；AES-192和AES-256则需要两个阶段联合工作。在方法论层面，本章详细阐述排除法的数学原理——从单样本分析到多样本统计收敛的完整推导链，以及等效块解密如何将缓存访问信息转化为倒数第二轮密钥约束条件。在系统实现层面，本章描述基于父子进程模型的系统架构、进程间同步通信协议、量化评估指标体系的设计与实现，以及CPU核心绑定、随机噪声注入和缓存刷新等缓解措施的工程实现方案。

\section{第一阶段：恢复最后一轮子密钥}

\subsection{AES最后一轮特征}

AES加密的最后一轮没有MixColumns。设最后一轮输入状态为$P$，密钥为$K_{last}$，输出密文为$C$：

\begin{equation}
C = ShiftRows(SubBytes(P)) \oplus K_{last}
\end{equation}

ShiftRows实现行位移，但T表可以将ShiftRows融合到查表映射中（详见下文）。设$S$为S盒，$P_i$表示与密文字节$C_i$对应的SubBytes输入字节（即经ShiftRows重排后产生$C_i$的那个输入字节），$K_{last,i}$为最后一轮密钥的第$i$个字节：

\begin{equation}
C_i = S(P_i) \oplus K_{last,i}
\end{equation}

\subsection{T表结构与监控策略}

T表加速用4张预计算表$Te_0$，$Te_1$，$Te_2$，$Te_3$。每张$Te$表256项×4字节=1024字节=16个缓存行。索引0-15占第0号缓存行，16-31占第1号，以此类推。攻击方仅监听每张T表的第0号缓存行，共4个。该信息量足以排除错误候选，同时可降低clflush开销和测量噪声。

仅监控第0缓存行的充分性分析，可以从两个角度：信息量和效率。

从信息量上看，每监测到一次未命中，就可以排除64个候选值，也就是排除一次$1/16$的搜索空间。为什么呢？以Te0为例：在Flush+Reload攻击中，如果Te0的缓存行未被访问，则意味着本次加密中所有使用Te0的输入字节都不在0～15范围内，而$Te0$负责的是第2行的输出，因此会影响第2行的4个密钥字节（即字节位置2，6，10，14）。同理：$Te1$未命中会影响第3行（字节3，7，11，15），$Te2$未命中会影响第0行（字节0，4，8，12），$Te3$未命中会影响第1行（字节1，5，9，13）。

需要特别说明，上述$Te_k$与状态行的对应关系是针对最后一轮的，与标准轮中$Te_k$对应第$k$行的规则不同。原因在于：本实验的受害者实现在最后一轮复用$Te$表并通过掩码提取S盒字节（详见3.1.3节），而每个$Te$表的S盒值位于32位条目的不同字节位置——$Te_0$的S盒值在字节1和2，$Te_2$的S盒值在字节0和3。最后一轮从$Te_2$提取字节0（最高有效字节）作为第0行输出，从$Te_3$提取字节1作为第1行输出，从$Te_0$提取字节2作为第2行输出，从$Te_1$提取字节3（最低有效字节）作为第3行输出。因此最后一轮的映射为$Te_0 \to$第2行、$Te_1 \to$第3行、$Te_2 \to$第0行、$Te_3 \to$第1行，恰好是标准轮映射的循环移位。排除法针对的是最后一轮密钥恢复，因此使用最后一轮的映射关系。第二阶段恢复倒数第二轮密钥时，目标轮为标准轮，映射回到$Te_k \to$第$k$行（详见3.2.2节）。

因此若$Te0$的缓存行未被访问就会影响第2行的4个密钥字节，对于给定的密文字节$C_i$，索引落在0~15对应的密钥候选值为$C_i \oplus \text{Sbox}[l]$（$l=0\dots15$），恰好16个候选。由于索引不可能为0~15，这16个候选值可在本次样本中被“排除”（即排除计数+1）。一次未命中的总排除量就等于4个字节$\times$16个候选=64个候选值。

如果监控全部16行会怎样？每次加密中4张T表需被访问多次，实际未被访问的缓存行数量较少。以具体数值为例：AES-128中$Te_0$在10轮中共被访问40次，每次独立以$1/16$概率落在任意给定缓存行。单个缓存行未被访问的概率$(15/16)^{40} \approx 0.076$，期望未命中缓存行数约$16 \times 0.076 \approx 1.2$。监控全部16行平均可获取约1.2个未命中缓存行的排除信息，而仅监控第0行平均可获取约0.076个。全监控的信息量约为单行监控的16倍，但监控成本亦为16倍——每单位成本的信息效率几乎相同。

从效率上看，监控$k$个缓存行意味着每次Flush+Reload需执行$k$次clflush和$k$次内存访问测量。clflush是串行化指令，延迟约100-200个CPU周期。监控4行（每张T表1行）的总刷新开销约400-800周期。监控64行（每张T表16行）则需6400-12800周期。后者不仅显著增大每次采样的时间成本，测量窗口亦随之扩大，更多噪声和抢占干扰将被引入。因此仅监控第0缓存行是信息量与测量效率之间的最优平衡点。

\subsection{最后一轮的T表访问方式与攻击前提}

标准AES最后一轮不含MixColumns，主流的高性能软件实现（如OpenSSL）通常在最后一轮使用一个独立的、仅包含S盒替换值的查找表（例如一个256字节的S盒），而非复用Te表——因为最后一轮没有MixColumns操作，使用Te表再提取S盒字节是多余的。但在本实验中的受害者实现并非如此：最后一轮仍访问$Te$表，通过位掩码仅提取每个32位条目中对应纯S盒值的那个字节。

该实现细节对Flush+Reload攻击是有利的：最后一轮虽无MixColumns，但仍会在$Te$表上留下缓存访问痕迹。监听方监控$Te$表即可捕获最后一轮的查表行为。而且由于最终仅使用S盒字节，密文字节与查表索引的关系不变：

\begin{equation}
C_i = S(P_i) \oplus K_{last,i}
\end{equation}

需要明确的是，本文的攻击框架针对的是一种在最后一轮仍复用$Te$表的AES实现。然而，攻击方法学本身是通用的：无论受害者最后一轮用$Te$表还是独立的$S$表，攻击者监控的目标内存区域和排除法的数学原理都同样适用，只需调整监控的查找表地址即可。如果实现改用独立S盒（比如OpenSSL），监听方将监控目标从$Te$表缓存行替换为S盒缓存行，排除公式依然成立。

\subsection{T表查表索引与密钥的代数关系}

要理解排除法为何能恢复密钥，需要从T表查表的代数结构出发，阐明密钥、明文和T表索引之间的推导链。

标准轮的T表查表过程。
AES标准轮（非最后一轮）对状态矩阵的每一列执行SubBytes、ShiftRows和MixColumns三个操作。以第$r$轮的第$j$列为例，设输入状态字节为$p_{0,j}$，$p_{1,j}$，$p_{2,j}$，$p_{3,j}$，轮密钥字节为$k_{0,j}$，$k_{1,j}$，$k_{2,j}$，$k_{3,j}$，T表查表过程为：

\begin{equation}
\begin{pmatrix} s'_{0,j} \\ s'_{1,j} \\ s'_{2,j} \\ s'_{3,j} \end{pmatrix} = Te_0[p_{0,j}] \oplus Te_1[p_{1,j+1}] \oplus Te_2[p_{2,j+2}] \oplus Te_3[p_{3,j+3}] \oplus \begin{pmatrix} k_{0,j} \\ k_{1,j} \\ k_{2,j} \\ k_{3,j} \end{pmatrix}
\end{equation}

其中$p_{i,j}$的下标$j$经过ShiftRows变换，列号偏移量分别为0、1、2、3。每个$Te_i$的查表索引就是对应的状态字节$p_{i,j+i}$。

最后一轮的T表查表索引。
如上节所述，受害者在最后一轮访问$Te$表并通过掩码提取S盒字节，因此查表索引就是状态字节$P_i$。设最后一轮输入为$P$，密文为$C$，最后一轮密钥为$K_{last}$：

\begin{equation}
C_i = S(P_i) \oplus K_{last,i}
\end{equation}

反解出最后一轮的查表索引：

\begin{equation}
P_i = S^{-1}(C_i \oplus K_{last,i})
\end{equation}

攻击者已知密文$C_i$，如果猜测到密钥候选$k$，即可计算对应的查表索引$S^{-1}(C_i \oplus k)$。该索引落在哪个缓存行，完全由索引的高4位决定：

\begin{equation}
CacheLine = \lfloor S^{-1}(C_i \oplus k) / 16 \rfloor
\end{equation}

这是因为T表的结构与缓存行的映射关系所决定的。每个T表有256个条目（索引0x00$\sim$0xFF），每个条目占4字节，总大小为1024字节。缓存行大小为64字节，每个缓存行可存放$64/4=16$个连续的T表条目，因此第0个缓存行存放索引0$\sim$15，第1个缓存行存放索引16$\sim$31，依此类推。一个8位的索引值（0$\sim$255）中，低4位（bit0$\sim$bit3）决定该条目在缓存行内的偏移位置，高4位（bit4$\sim$bit7）决定该条目属于第几个缓存行。因此缓存行号等于索引右移4位，即$\lfloor \text{索引}/16 \rfloor$。在攻击中，攻击者无法直接获知完整的8位索引，但通过Flush+Reload观察到某个缓存行是否被访问，就能推断索引的高4位信息。例如，若第0号缓存行未被访问，则索引不可能是0$\sim$15，从而排除16个候选值——这正是排除法能够逐步收敛密钥的根本原因。

为方便后续表述，将上面两个公式写为：

\begin{align}
\text{index}_k &= S^{-1}(C_i \oplus k) \\
H(\text{index}_k) &= \lfloor S^{-1}(C_i \oplus k) / 16 \rfloor
\end{align}

通过Flush+Reload观测该T表第0缓存行的访问状态（被访问即命中，未被访问即未命中）。若第0缓存行未被访问，则实际加密过程中没有任何索引落在该行，所有$H(S^{-1}(C_i \oplus k)) = 0$的候选$k$与观测矛盾，予以排除；若被访问，则无法确定具体是哪个索引触发了访问，不提供排除信息。

排除条件仅在未命中时生效：

\begin{equation}
\text{排除候选} \Leftrightarrow \text{该T表第0缓存行未被访问} \land H(S^{-1}(C_i \oplus k)) = 0
\end{equation}

\subsection{排除法的数学推导}

在多样本累积与统计收敛中，为每个密钥字节初始化一个排除计数器数组，256个候选初始全0。

对于$N$个样本，第$j$个样本的密文字节是$C_i^{(j)}$。$m^{(j)}$表示第$j$次加密中该T表第0缓存行是否没被访问（$m^{(j)} = 1$是未命中，$m^{(j)} = 0$是命中）。排除计数$CK[i][k]$定义为密钥字节$i$的候选$k$被排除的总次数：

\begin{equation}
CK[i][k] = \sum_{j=1}^{N} m^{(j)} \cdot \mathbf{1}[H(S^{-1}(C_i^{(j)} \oplus k)) = 0]
\end{equation}

$\mathbf{1}[\cdot]$是指示函数，条件成立取1否则取0。命中时$m^{(j)} = 0$，整项为0，不计数。未命中时$m^{(j)} = 1$，若$H(S^{-1}(C_i^{(j)} \oplus k)) = 0$则计数加1。设$k^*$为密钥字节$i$的真实值，正确候选$k^*$始终不会被排除。其原因在于：正确索引$S^{-1}(C_i^{(j)} \oplus k^*)$对应的缓存行一定被访问过，未命中时$H(S^{-1}(C_i^{(j)} \oplus k^*)) \neq 0$，因此$CK[i][k^*] = 0$。错误候选的$CK[i][k] > 0$。计数最少的候选就是要恢复的密钥字节。

期望样本量分析。
排除法仅在未命中样本上执行排除，命中样本直接略过。设总样本量$N$，单次加密中某T表第0缓存行不被访问的概率记作$p_{miss}$。以下推导基于查表索引均匀分布的假设——在攻击者已知密文、密钥未知且服从均匀分布的条件下，该假设基本成立。单次T表查表索引落在第0缓存行的概率是$1/16$，不落的是$15/16$。

由于每次加密中同一$Te$表被多个字节多次访问（每轮4次，4列各1次），AES-128共有9轮标准轮（含MixColumns）加1轮最终轮（无MixColumns），本实现最终轮仍复用$Te$表，故每张$Te$表被访问次数为：标准轮$9 \times 4 = 36$次，最终轮4次，合计40次；AES-192为11标准轮加1最终轮，每表$11 \times 4 + 4 = 48$次；AES-256为13标准轮加1最终轮，每表$13 \times 4 + 4 = 56$次。在均匀随机假设下，单次访问不落在第0缓存行的概率为$15/16$，因此$p_{miss} = (15/16)^n$，其中$n$为该$Te$表的总访问次数。代入得AES-128的$p_{miss} \approx 0.076$，AES-192的$p_{miss} \approx 0.045$，AES-256的$p_{miss} \approx 0.027$。期望有效样本数约为$N \cdot p_{miss}$。需要注意，上述$n$的取值以受害者实现在所有轮次（包括最后一轮）均访问$Te$表为前提；若最后一轮改用独立S盒表，则$Te$表的总访问次数减少4次（最后一轮不再查$Te$表），$p_{miss}$将略有增大，但推导方法不变。

每个有效样本对单个密钥字节可排除$16/256 = 1/16$的候选。对某个错误候选$k$，单次样本将其排除的概率为$p_{miss}/16$（需同时满足未命中且$H(S^{-1}(C_i \oplus k)) = 0$），保留概率为$1 - p_{miss}/16$。初始候选集256个，经过$N$个样本后期望剩余的候选数：

\begin{equation}
E[\text{remaining}] = 256 \cdot \left(1 - \frac{p_{miss}}{16}\right)^{N}
\end{equation}

要让剩余候选数等于1，代入上述$p_{miss}$估计，由$E[\text{remaining}] = 256 \cdot (1 - p_{miss}/16)^N = 1$解得$N \approx \ln(1/256) / \ln(1 - p_{miss}/16)$，当$p_{miss}$较小时$\ln(1 - p_{miss}/16) \approx -p_{miss}/16$，线性近似为$N \approx 16\ln(256) / p_{miss} \approx 89 / p_{miss}$。AES-128需总样本量约$89 / 0.076 \approx 1171$，AES-192约$89 / 0.045 \approx 1978$，AES-256约$89 / 0.027 \approx 3296$。

需要强调，上述$89/p_{miss}$仅是单字节独立恢复的近似下界，不能直接作为整体攻击的样本量推荐。整体攻击需要16个密钥字节同时恢复，联合成功概率远低于单字节概率。假设单字节恢复成功率为$p_{byte}$，整体成功率约为$p_{byte}^{16}$；对于两阶段攻击（AES-192/256），整体成功率进一步约为$p_{byte}^{32}$。因此实际需要的样本量远超上述估算，原因包括多个层面：样本之间存在相关性、测量本身存在噪声、需要保留安全冗余，其中最重要的是16个字节需同时恢复，任意字节出错即导致整体失败，因此需要远多于单字节理论值的样本数以提升联合置信度。本实验里AES-128、AES-192、AES-256分别使用5000、9000、15000个样本。

经过大量样本累积后，理论上真实值$k^*$的排除计数始终为0，而错误候选的排除计数逐渐增大。最终排除计数最少的候选即为恢复的密钥字节。

误判率分析。
排除法的可靠性取决于两个条件：正确候选始终不被排除，错误候选被充分排除。实际测量中存在噪声，可能导致两种误判：

假阴性：正确候选$k^*$被错误排除。根源在于$Te$表第0缓存行实际上被访问了，但监听方测到了未命中。Flush和Reload之间的时间窗口内，缓存行可能被其他进程驱逐、CPU预取器换出，或者OS调度导致Victim尚未完成加密，监听方即开始Reload。假阴性率记作$P_{FN}$，它使$CK[i][k^*] > 0$，正确候选的排除计数不再为零。

假阳性：错误候选$k$未被排除。这发生在$Te$表第0缓存行实际上未被访问，但监听方测到了命中。原因有多种：其他进程（如共享库代码）恰好访问了同一缓存行、测量噪声使访问时间低于阈值。假阳性率记作$P_{FP}$，它使错误候选的排除计数偏低，收敛速度减慢。假阳性不直接导致错误恢复，但需累积更多样本方能将错误候选与正确候选区分开。

设单次测量的假阴性率$P_{FN}$、假阳性率$P_{FP}$，有效未命中概率修正为$p'_{miss} = p_{miss}(1 - P_{FP}) + (1 - p_{miss})P_{FN}$。低噪声环境（$P_{FN} < 0.01$，$P_{FP} < 0.01$）下，修正量对收敛速度的影响可以忽略。高噪声环境则需按比例增加样本量予以补偿。本实验在可控环境中进行，噪声水平低，假阴性和假阳性率均控制在1\%以下。

\subsection{AES-128密钥恢复}

AES-128恢复最后一轮密钥$K_{10}$以后，通过密钥扩展的逆过程即可获得原始密钥$K_0$。AES-128密钥扩展生成44个32位字$w_0$，$w_1$，$...$，$w_{43}$，$K_{10}$对应$w_{40}$，$w_{41}$，$w_{42}$，$w_{43}$。

密钥扩展逆过程：

\begin{equation}
w_{i-4} = \begin{cases} w_i \oplus SubWord(RotWord(w_{i-1})) \oplus Rcon[i/4] & \text{if } i \mod 4 = 0 \\ w_i \oplus w_{i-1} & \text{otherwise} \end{cases}
\end{equation}

其中$Rcon[r]$为AES第$r$轮的轮常量，$SubWord$对32位字的每个字节做S盒替换，$RotWord$将32位字循环左移一个字节。

从$w_{43}$逐步逆向推到$w_0$，$w_1$，$w_2$，$w_3$，即可获得原始密钥$K_0$。

\section{第二阶段：恢复倒数第二轮子密钥}

AES-192和AES-256仅恢复最后一轮子密钥无法反推原始密钥，还需获取倒数第二轮子密钥。攻击者使用同一批Flush+Reload监控数据，结合已知的最后一轮密钥和密文，计算解密中间值，进而恢复倒数第二轮密钥。

\subsection{复用监控信息与等效块解密}

第一阶段Flush+Reload仅记录了整次加密中各T表缓存行是否被访问，无法区分各轮次的访问。攻击者无法从中提取倒数第二轮独有的访问特征。然而从另一角度考虑：缓存行未被访问，即意味着所有轮次均未访问该行。攻击者可将该未命中信息应用于倒数第二轮的中间状态，通过排除法恢复倒数第二轮密钥。

还有一个隐含前提：受害者AES实现在所有轮次中访问的是同一组$Te$表（标准轮使用完整的32位查表结果，最后一轮通过掩码仅提取S盒字节）。这样才能保证Flush+Reload监控到的$Te$表缓存行未命中信息对所有轮次均有效。

等效块解密方法的推导思路。
第二阶段的核心方法为等效块解密，攻击者无需执行完整AES解密，仅需计算密文到倒数第二轮加密输入状态$S_{r-1}$的映射即可。$S_{r-1}$的每个字节，恰好就是第$r-1$轮查$Te$表时的输入索引。

推导的关键在于理解AES等效解密结构。标准AES解密依次执行InvShiftRows、InvSubBytes、AddRoundKey、InvMixColumns。但该顺序与加密过程的逆序并非直接对应：加密中SubBytes在ShiftRows之前，解密中InvShiftRows在InvSubBytes之前。为使解密结构更规整，NIST标准定义了"等效解密"：将InvSubBytes和InvShiftRows交换位置（二者可交换，因为InvSubBytes逐字节操作，InvShiftRows仅改变位置），然后将轮密钥做InvMixColumns变换，得到$Kd_r = InvMixColumns(K_r)$。

在等效解密框架下，从密文$C$出发，经过一轮等效解密得到的中间值$D_2$恰好等于$ShiftRows(SubBytes(S_{r-1}))$。这正是第$r-1$轮加密中SubBytes和ShiftRows的输出。由此可以直接反推$S_{r-1}$，而$S_{r-1}$的每个字节就是第$r-1$轮查$Te$表的输入索引。攻击者只需已知密文$C$和最后一轮密钥$K_{last}$，加上猜测的倒数第二轮密钥$K_{last-1}$（或其等效形式$Kd_1$），即可在本地完成所有计算。

AES-128为何不需要第二阶段。AES-128的密钥扩展生成44个32位字$w_0$，$w_1$，$...$，$w_{43}$，原始密钥$K_0$对应$w_0$，$w_1$，$w_2$，$w_3$，最后一轮密钥$K_{10}$对应$w_{40}$，$w_{41}$，$w_{42}$，$w_{43}$。密钥扩展的递推关系为每4个字一组，$w_i$仅依赖于$w_{i-4}$和$w_{i-1}$。从$w_{40}$，$w_{41}$，$w_{42}$，$w_{43}$出发，可以逐步逆向推导出$w_{36}$，$w_{32}$，$...$，$w_0$，完整恢复全部44个字。因此AES-128仅凭最后一轮密钥$K_{10}$即可唯一确定原始密钥$K_0$，无需恢复倒数第二轮密钥。

AES-192和AES-256则不同——仅凭最后一轮子密钥的4个字，无法逆向推出完整原始密钥。这是由每个AES变体的密钥长度$Nk$与每轮子密钥长度（固定4字）之间的不匹配造成的。密钥扩展逆过程需要$Nk$个连续字才能回溯，但最后一轮子密钥只提供4个字，当$Nk > 4$时就会出现缺口。这就是AES-192/256必须额外恢复倒数第二轮子密钥的根本原因。下面结合逆推公式做具体说明。

AES-192的密钥扩展。$Nk = 6$，一共生成52个字$w_0, \dots, w_{51}$。递推关系以6个字为一组：

\begin{equation}
w_i = \begin{cases} w_{i-6} \oplus SubWord(RotWord(w_{i-1})) \oplus Rcon[i/6] & i \bmod 6 = 0 \\ w_{i-6} \oplus w_{i-1} & \text{否则} \end{cases}
\end{equation}

逆过程只需将等式移项：

\begin{equation}
w_{i-6} = \begin{cases} w_i \oplus SubWord(RotWord(w_{i-1})) \oplus Rcon[i/6] & i \bmod 6 = 0 \\ w_i \oplus w_{i-1} & \text{否则} \end{cases}
\end{equation}

攻击者从$K_{12} = (w_{48}, w_{49}, w_{50}, w_{51})$出发，共有4个已知字。逆推过程如图~\ref{fig:aes192_reverse}所示：

\begin{figure}[htbp]
\centering
\begin{tikzpicture}[node distance=0.8cm, auto,
    block/.style={rectangle, draw, minimum width=1.8cm, minimum height=0.7cm, font=\small},
    known/.style={fill=green!20},
    partial/.style={fill=yellow!20},
    unknown/.style={fill=red!10}]

% w indices
\node[block, known] (w51) {$w_{51}$ (已知)};
\node[block, known, right=of w51] (w50) {$w_{50}$ (已知)};
\node[block, known, right=of w50] (w49) {$w_{49}$ (已知)};
\node[block, known, right=of w49] (w48) {$w_{48}$ (已知)};

\node[block, partial, below=of w51] (w45) {$w_{45}$ 可求};
\node[block, partial, below=of w50] (w44) {$w_{44}$ 可求};
\node[block, partial, below=of w49] (w43) {$w_{43}$ 可求};
\node[block, unknown, below=of w48] (w42) {$w_{42}$ 不可求};

\node[block, unknown, below=of w45] (w46) {$w_{46}$ 不可求};
\node[block, unknown, below=of w44] (w47) {$w_{47}$ 不可求};

\draw[->] (w51) -- (w45) node[midway, right] {\tiny $w_{51}\!\oplus\!w_{50}$};
\draw[->] (w50) -- (w44) node[midway, right] {\tiny $w_{50}\!\oplus\!w_{49}$};
\draw[->] (w49) -- (w43) node[midway, right] {\tiny $w_{49}\!\oplus\!w_{48}$};
\draw[->] (w48) -- (w42) node[midway, right] {\tiny 需$w_{47}$};

\end{tikzpicture}
\caption{AES-192从$K_{12}$的4个已知字逆向推导。$w_{45}, w_{44}, w_{43}$可直接异或求得；$w_{42}$需要$w_{47}$但$w_{47}$未知，逆推中断。$w_{46}, w_{47}$也无法从已有信息获取}
\label{fig:aes192_reverse}
\end{figure}

具体地：
\begin{itemize}
    \item $w_{51} = w_{45} \oplus w_{50} \;\Rightarrow\; w_{45} = w_{51} \oplus w_{50}$，可求。
    \item $w_{50} = w_{44} \oplus w_{49} \;\Rightarrow\; w_{44} = w_{50} \oplus w_{49}$，可求。
    \item $w_{49} = w_{43} \oplus w_{48} \;\Rightarrow\; w_{43} = w_{49} \oplus w_{48}$，可求。
    \item $w_{48} = w_{42} \oplus SubWord(RotWord(w_{47})) \oplus Rcon[8]$：$w_{47}$未知，$\mathbf{w_{42}}$不可求。
\end{itemize}
到此逆推中断——$w_{42}$求不出，$w_{46}$和$w_{47}$也无从获得，无法继续回到$w_0 \sim w_5$。唯有同时恢复$K_{11} = (w_{44}, w_{45}, w_{46}, w_{47})$，凑齐$w_{44} \sim w_{51}$共8个连续字，方可沿逆过程一路回推到$w_0$。

AES-256的密钥扩展。$Nk = 8$，一共生成60个字$w_0, \dots, w_{59}$。AES-256多了$i \bmod 8 = 4$这一特例（需经过$SubWord$而非裸异或）。递推关系：

\begin{equation}
w_i = \begin{cases} w_{i-8} \oplus SubWord(RotWord(w_{i-1})) \oplus Rcon[i/8] & i \bmod 8 = 0 \\ w_{i-8} \oplus SubWord(w_{i-1}) & i \bmod 8 = 4 \\ w_{i-8} \oplus w_{i-1} & \text{否则} \end{cases}
\end{equation}

逆过程：

\begin{equation}
w_{i-8} = \begin{cases} w_i \oplus SubWord(RotWord(w_{i-1})) \oplus Rcon[i/8] & i \bmod 8 = 0 \\ w_i \oplus SubWord(w_{i-1}) & i \bmod 8 = 4 \\ w_i \oplus w_{i-1} & \text{否则} \end{cases}
\end{equation}

攻击者从$K_{14} = (w_{56}, w_{57}, w_{58}, w_{59})$出发：
\begin{itemize}
    \item $w_{59} = w_{51} \oplus w_{58} \;\Rightarrow\; w_{51}$可求。
    \item $w_{58} = w_{50} \oplus w_{57} \;\Rightarrow\; w_{50}$可求。
    \item $w_{57} = w_{49} \oplus w_{56} \;\Rightarrow\; w_{49}$可求。
    \item $w_{56}$需要$w_{48}$和$w_{55}$，$w_{55}$未知，$\mathbf{w_{48}}$不可求。
\end{itemize}
同理，$w_{52}, w_{53}, w_{54}, w_{55}$均无法从4个已知字中推出，逆推在$w_{48}$处中断。必须同时恢复$K_{13} = (w_{52}, w_{53}, w_{54}, w_{55})$，获得$w_{52} \sim w_{59}$共8个连续字才能完整逆推。

小结。AES-128的$Nk = 4$，一轮子密钥的4个字恰好等于$Nk$，逆推通路。AES-192的$Nk = 6$、AES-256的$Nk = 8$，$Nk > 4$导致一轮子密钥信息不足，必须额外恢复倒数第二轮子密钥，补齐$Nk$个连续字才能打通逆推链路。这正是两阶段攻击的数学根源。

\subsection{从密文到$S_{13}$的推导与缓存行映射}

以AES-256为例。攻击者在第一阶段获得了密文$C$和最后一轮子密钥$K_{14}$，需要从密文推导出倒数第二轮加密的输入状态$S_{13}$，并建立$S_{13}$到Flush+Reload可观测缓存行的映射关系。

第一步：密文与最后一轮密钥异或。
\begin{equation}
D_{1,i} = C_i \oplus K_{14,i}, \quad i = 0, 1, \dots, 15
\end{equation}
攻击者已知$C$和$K_{14}$，$D_1$的每个字节均可精确计算。

第二步：通过逆T表计算中间值$IV_1$。
AES解密使用4张逆T表$Td_0, Td_1, Td_2, Td_3$，每张将逆S盒替换与InvMixColumns系数融合为单次查表，结构与加密T表$Te$对称。$IV_1$的每个字节由4次逆T表查表异或得到：

\begin{align}
IV_{1,i} &= Td_{\sigma(i,0)}[D_1[\pi(i,0)]] \oplus Td_{\sigma(i,1)}[D_1[\pi(i,1)]] \nonumber \\
&\quad \oplus Td_{\sigma(i,2)}[D_1[\pi(i,2)]] \oplus Td_{\sigma(i,3)}[D_1[\pi(i,3)]]
\end{align}

其中$\sigma(i,\cdot)$和$\pi(i,\cdot)$分别由InvShiftRows和逆列混合的矩阵结构决定。以第0字节为例：

\begin{equation}
IV_{1,0} = Td_0[D_1[0]] \oplus Td_1[D_1[13]] \oplus Td_2[D_1[10]] \oplus Td_3[D_1[7]]
\end{equation}

此处$Td_k$的下标对应InvShiftRows后第$k$行的字节。需要指出，部分实现（如将InvShiftRows吸收到$Td$表定义内部的方案）中$Td_1$和$Td_3$的系数行互换，此时$D_1$的索引变为对称形式，功能上完全等价。本文沿用$Td_k$下标与InvShiftRows后行号一致的约定。

第三步：异或等效解密密钥得到$D_2$。
\begin{equation}
D_{2,i} = IV_{1,i} \oplus Kd_{1,i}
\end{equation}
其中等效解密轮密钥$Kd_1 = InvMixColumns(K_{13})$，即对$K_{13}$的每一列施加InvMixColumns矩阵变换：

\begin{equation}
\begin{pmatrix} Kd_{1,0,j} \\ Kd_{1,1,j} \\ Kd_{1,2,j} \\ Kd_{1,3,j} \end{pmatrix} =
\begin{pmatrix} \mathtt{0x0E} & \mathtt{0x0B} & \mathtt{0x0D} & \mathtt{0x09} \\ \mathtt{0x09} & \mathtt{0x0E} & \mathtt{0x0B} & \mathtt{0x0D} \\ \mathtt{0x0D} & \mathtt{0x09} & \mathtt{0x0E} & \mathtt{0x0B} \\ \mathtt{0x0B} & \mathtt{0x0D} & \mathtt{0x09} & \mathtt{0x0E} \end{pmatrix}
\begin{pmatrix} K_{13,0,j} \\ K_{13,1,j} \\ K_{13,2,j} \\ K_{13,3,j} \end{pmatrix}
\end{equation}

第四步：从$D_2$反推$S_{13}$。
加密第13轮依次为$SubBytes \rightarrow ShiftRows \rightarrow MixColumns \rightarrow AddRoundKey(K_{13})$。由等效解密的结构关系：
\begin{equation}
D_2 = ShiftRows(SubBytes(S_{13}))
\end{equation}
逐字节反推：
\begin{equation}
S_{13,i} = S^{-1}(D_{2, \rho(i)})
\end{equation}
其中$\rho(i)$为InvShiftRows的字节重排映射。$S_{13}$的每个字节就是第13轮加密时查$Te$表的输入索引。

第五步：$S_{13}$到T表缓存行的映射。
第13轮为标准轮，状态矩阵按行划分，每行字节查同一张$Te$表：
\begin{equation}
S_{13,i} \rightarrow Te_{i \bmod 4}
\end{equation}
即第0行字节（$i = 0, 4, 8, 12$）查$Te_0$，第1行字节（$i = 1, 5, 9, 13$）查$Te_1$，第2行字节（$i = 2, 6, 10, 14$）查$Te_2$，第3行字节（$i = 3, 7, 11, 15$）查$Te_3$。索引的高4位决定缓存行：
\begin{equation}
H(S_{13,i}) = S_{13,i} \gg 4
\end{equation}
攻击者无法直接观测$S_{13,i}$，但可通过Flush+Reload的缓存行未命中信息间接排除：若$Te_{i \bmod 4}$的第0缓存行在整次加密中全局未被访问，则第13轮中$H(S_{13,i}) \neq 0$。

\subsection{复用监控数据直接恢复$Kd_1$}

第一阶段Flush+Reload记录了加密中每次对$Te$表的缓存访问。某个$Te$表的第0缓存行，如果在整次加密中均未被访问过，则不仅是最后一轮，所有轮次（包括第13轮）均未访问该行。

由3.2.2节的推导，$Kd_{1,i} = IV_{1,i} \oplus D_{2,i}$，而$D_{2,i} = S(S_{13,\rho^{-1}(i)})$（经ShiftRows重排后的S盒输出）。因此，枚举$S_{13}$的候选值等价于枚举$Kd_1$的候选值。本实现采用直接恢复$Kd_1$的策略：对$Kd_1$的每个字节枚举256个候选值，利用缓存行未命中信息排除不可能的候选，无需先恢复$S_{13}$再推导$Kd_1$。

以$Kd_1$的某个字节为例说明排除过程。该字节位于第$row$行、第$col$列，对应状态矩阵中位置$i = row + 4 \cdot col$。在第$j$次加密中，攻击者已知$IV_{1,i}^{(j)}$和该次Flush+Reload的观测结果。关键推导分三步：

从缓存观测推导$S_{13,i}$的约束。$S_{13,i}$是第13轮查$Te_{row}$表的输入索引。若$Te_{row}$的第0缓存行在本次加密中全局未命中，则所有轮次（包括第13轮）均未访问该缓存行，即第$row$行所有字节的查表索引均不落在0$\sim$15范围内。因此$S_{13,i} \neq l$（$l = 0, \dots, 15$）。

从$S_{13,i}$的约束推导$D_{2,\rho(i)}$的约束。$D_{2,\rho(i)} = S(S_{13,i})$（S盒替换，索引由ShiftRows重排）。由$S_{13,i} \neq l$得出$D_{2,\rho(i)} \neq \text{Sbox}[l]$（$l = 0, \dots, 15$）。

从$D_{2,\rho(i)}$的约束推导$Kd_{1,i}$的排除。由3.2.2节的推导，$Kd_{1,i} = IV_{1,i} \oplus D_{2,\rho(i)}$。因此$Kd_{1,i}$不能等于$IV_{1,i} \oplus \text{Sbox}[l]$，这16个候选值可在本次样本中被排除。

定义$Kd_1$的排除计数$CKd_1[i][k]$为第$i$字节的候选值$k$被排除的累积次数：
\begin{equation}
CKd_1[i][k] = \sum_{j=1}^{N} m_{row}^{(j)} \cdot \mathbf{1}\bigl[k \in \{IV_{1,i}^{(j)} \oplus \text{Sbox}[l] \mid l = 0, \dots, 15\}\bigr]
\end{equation}
其中$m_{row}^{(j)}$表示第$j$次加密中$Te_{row}$的第0缓存行是否未被访问（1为未命中，0为命中），$\mathbf{1}[\cdot]$为指示函数。正确候选的排除计数趋近于零，经过足够多样本累积后，排除计数最少的候选即为恢复的$Kd_{1,i}$值。

上述直接恢复$Kd_1$的方法与"先恢复$S_{13}$再推导$Kd_1$"在数学上完全等价：排除$S_{13,i}$的候选索引$s$（满足$H(s)=0$），等价于排除$Kd_1$候选$IV_{1,i} \oplus \text{Sbox}[s]$。直接法省去了中间步骤，实现更简洁。

\subsection{从$Kd_1$计算倒数第二轮子密钥}

$Kd_1$恢复出来后，对其每一列做MixColumns变换即可得到加密倒数第二轮子密钥$K_{13}$（AES-256）或$K_{11}$（AES-192）：

写成矩阵形式，对$Kd_1$的每一列做MixColumns：

\begin{equation}
\begin{pmatrix} K_{13,0,j} \\ K_{13,1,j} \\ K_{13,2,j} \\ K_{13,3,j} \end{pmatrix} = \begin{pmatrix} 0x02 & 0x03 & 0x01 & 0x01 \\ 0x01 & 0x02 & 0x03 & 0x01 \\ 0x01 & 0x01 & 0x02 & 0x03 \\ 0x03 & 0x01 & 0x01 & 0x02 \end{pmatrix} \begin{pmatrix} Kd_{1,0,j} \\ Kd_{1,1,j} \\ Kd_{1,2,j} \\ Kd_{1,3,j} \end{pmatrix}
\end{equation}

\subsection{从轮密钥恢复原始密钥}

AES-192恢复$K_{12}$和$K_{11}$后获得连续8个32位字（$w_{44}$--$w_{51}$），足以完整逆推192位原始密钥$K_0$；AES-256恢复$K_{14}$和$K_{13}$后获得连续8个32位字（$w_{52}$--$w_{59}$），同样可完整逆推256位原始密钥$K_0$。

AES-192逆向推导：

\begin{equation}
w_{i-6} = \begin{cases} w_i \oplus SubWord(RotWord(w_{i-1})) \oplus Rcon[i/6] & \text{if } i \mod 6 = 0 \\ w_i \oplus w_{i-1} & \text{otherwise} \end{cases}
\end{equation}

AES-256逆向推导：

\begin{equation}
w_{i-8} = \begin{cases} w_i \oplus SubWord(RotWord(w_{i-1})) \oplus Rcon[i/8] & \text{if } i \mod 8 = 0 \\ w_i \oplus SubWord(w_{i-1}) & \text{if } i \mod 8 = 4 \\ w_i \oplus w_{i-1} & \text{otherwise} \end{cases}
\end{equation}

\subsection{两阶段攻击流程总结}

AES-128：使用5000样本，排除法恢复$K_{10}$，逆过程恢复$K_0$。

AES-192：使用9000样本，阶段一恢复$K_{12}$，阶段二复用监控数据直接恢复$Kd_1$，经MixColumns计算$K_{11}$，逆向恢复$K_0$。

AES-256：使用15000样本，阶段一恢复$K_{14}$，阶段二复用监控数据直接恢复$Kd_1$，经MixColumns计算$K_{13}$，逆向恢复$K_0$。

关键点在于：两个阶段共享同一批监控数据。攻击者只需执行一遍Flush+Reload，即可同时恢复两个轮密钥。

以AES-256为例，整个密钥恢复的数学过程可概括为以下公式链。设$C^{(j)}$为第$j$次加密的密文，$m_{k}^{(j)} \in \{0,1\}$表示$Te_k$的第0缓存行是否未被访问。

阶段一：恢复最后一轮子密钥$K_{14}$。
对每个密钥字节$i$（$i = 0,\dots,15$），利用最后一轮映射关系（3.1.3节），若对应$Te$表的第0缓存行未命中，则排除满足$H(S^{-1}(C_i^{(j)} \oplus k)) = 0$的候选$k$：
\begin{equation}
CK[i][k] = \sum_{j=1}^{N} m_{*}^{(j)} \cdot \mathbf{1}\bigl[H(S^{-1}(C_i^{(j)} \oplus k)) = 0\bigr]
\end{equation}
排除计数最少的候选即为恢复的$K_{14,i}$。

阶段二：直接恢复等效解密密钥$Kd_1$。
对$Kd_1$的每个字节$i$（$i = 0,\dots,15$），利用标准轮映射（$i$所在行$row$对应$Te_{row}$），若$Te_{row}$的第0缓存行未命中，则排除16个候选值$IV_{1,i}^{(j)} \oplus \text{Sbox}[l]$（$l = 0,\dots,15$）：
\begin{equation}
CKd_1[i][k] = \sum_{j=1}^{N} m_{row}^{(j)} \cdot \mathbf{1}\bigl[k \in \{IV_{1,i}^{(j)} \oplus \text{Sbox}[l] \mid l = 0,\dots,15\}\bigr]
\end{equation}

密钥恢复链。
\begin{equation}
Kd_1 \xrightarrow{\text{MixColumns}} K_{13}, \quad
\{K_{13}, K_{14}\} \xrightarrow{\text{密钥扩展逆过程}} K_0
\end{equation}

AES-192的结构与此完全一致，仅将$K_{14}$换为$K_{12}$、$K_{13}$换为$K_{11}$、$Kd_1$中的$K_{13}$换为$K_{11}$。AES-128则只需执行阶段一，由$K_{10}$直接逆推$K_0$。

\section{系统架构设计}

系统按模块划分为六个部分。共享内存模块创建和管理共享内存，放置T表数据。Victim进程模块模拟受害者，执行AES加密。Attacker进程模块执行Flush+Reload，收集缓存访问信息。密钥恢复模块从收集的信息中恢复AES密钥。量化评估模块计算各项指标，评估攻击效果。自动化实验模块支持批量测试和结果可视化。

系统采用父子进程模型，通过fork创建Victim进程。父进程为Attacker，负责刷新缓存、收集样本、恢复密钥。子进程为Victim，执行AES加密响应Attacker的请求。进程间通过管道同步通信，确保攻击流程的时序正确。

\subsection{进程间通信协议}

Attacker和Victim之间通过两对管道（pipe）实现同步，分别用于请求-响应方向。管道1（Attacker$\rightarrow$Victim）传递加密请求信号，管道2（Victim$\rightarrow$Attacker）传递加密完成通知信号。

单次采样的通信过程是：Attacker先执行Flush阶段，用clflush刷新目标T表缓存行；然后通过管道1向Victim发送1字节的请求信号（\texttt{`E'}表示"Encrypt"）；Victim从管道1读取请求信号后随即执行一次AES加密；加密完成后，Victim将16字节密文写入共享内存的指定字段，然后通过管道2向Attacker发送1字节的响应信号（\texttt{`D'}表示"Done"）；Attacker从管道2读取响应信号后，立即执行Reload阶段，测量目标缓存行的访问时间；然后从共享内存中读取密文，记录本次采样的缓存命中/未命中状态和密文，进入下一轮采样。密文通过共享内存传递而非管道，避免了管道传输16字节数据的额外延迟。

管道通信采用阻塞读写模式，天然实现了进程间的同步：Attacker在发送请求后阻塞等待Victim的响应，Victim在处理完前一个请求前不会接收新请求。这种设计避免了忙等待（busy-wait）带来的CPU资源浪费和时序干扰。管道缓冲区大小为64KB（Linux默认值），远大于单次通信的数据量，不会发生写阻塞。

\section{量化评估指标体系}

\subsection{基础指标}

样本总数sample\_count、测量次数measurement\_count（样本数×4）、命中次数hit\_count、未命中次数miss\_count、命中率hit\_rate（命中次数÷总测量次数）、未命中率miss\_rate。

样本总数反映攻击的数据需求，测量次数将样本数按监控目标数（4张T表）展开，反映总测量开销。命中率和未命中率直接决定排除法的收敛速度：未命中率越高，有效排除样本越多，收敛越快。在理想无噪声环境下，未命中率的理论值可由前文中的$p_{miss}$公式预测，实际测量值与理论值的偏差反映了噪声水平。

\subsection{时间统计指标}

基本统计量（均值、标准差、最小值、最大值）和分位数指标（中位数、Q1、Q3、P5、P95）。本实验使用分位数而不使用均值标准差，因为分位数对异常值的抵抗力强得多，更适合用户态高精度时间测量的场景。更进一步讲，用户态时间测量受操作系统调度、中断和动态频率调节的影响，异常值不可避免。均值和标准差对异常值敏感，少量极端值即可大幅偏移均值。分位数（特别是中位数和四分位距）对异常值具有鲁棒性，能更准确地反映典型访问时间的分布特征。P5和P95则提供了分布尾部的信息，有助于区分缓存命中和未命中的时间阈值选取。时间统计指标与基础指标之间存在关联：当命中率异常偏高时，访问时间的分布会向低值偏移（因为缓存命中时间远低于未命中时间），此时P95可以作为判断噪声水平的辅助指标。

\subsection{泄露带宽}

泄露带宽（Leakage Bandwidth）定义为单位时间内侧信道攻击从目标进程中提取的信息量，单位为比特/秒（bps）。它综合了两个因素：每次观测所获取的信息量（互信息）和单位时间内的观测次数（采样频率）。泄露带宽越高，攻击者在单位时间内获取的密钥信息越多，攻击速度越快。

\begin{equation}
BW = 4 \cdot MI_{binary} \cdot f_{sample}
\end{equation}

该公式的含义是：一次完整采样对4个缓存行（$Te_0$、$Te_1$、$Te_2$、$Te_3$各1个）各进行一次观测，每次观测提供$MI_{binary}$比特信息，故一次采样共获取$4 \cdot MI_{binary}$比特（假设各观测独立）；每秒完成$f_{sample}$次采样，两者相乘即为每秒获取的总信息量。因此，提升单次观测的互信息（降低噪声）或提升采样频率（缩短单次采样耗时）都能增大泄露带宽，加速密钥恢复。

\paragraph{互信息}

互信息用于量化单次Flush+Reload观测所获取的信息量。下面分步解释其含义与计算方法。

第一步：理解观测过程。每次Flush+Reload观测，攻击者执行以下操作：先用clflush将目标缓存行逐出缓存，然后等待Victim执行一次AES加密，最后测量该缓存行的访问时间。如果访问时间短（命中），说明Victim访问了该缓存行；如果访问时间长（未命中），说明Victim未访问。因此，单次观测的结果是一个二值信息：命中或未命中。

第二步：建立信息模型。将上述观测过程建模为二值信道。信道的输入是Victim的真实访问状态：设$X=1$表示Victim访问了目标缓存行，$X=0$表示未访问。信道的输出是攻击者的观测结果：设$Y=1$表示命中，$Y=0$表示未命中。攻击者的目标是通过观测$Y$来推断$X$，因为$X$与密钥相关——Victim访问哪个T表缓存行取决于密钥字节值，观测到"未命中"可以排除一个错误密钥候选。

第三步：确定信道参数。信道的特性由条件概率$P(Y|X)$描述。当Victim访问了目标缓存行（$X=1$）时，攻击者必然观测到命中（$Y=1$），即$P(Y=1|X=1)=1$。当Victim未访问（$X=0$）时，若无噪声则攻击者应观测到未命中，但噪声注入会以概率$f$产生假命中，即$P(Y=1|X=0)=f$（无噪声时$f=0$）。输入$X$的先验概率由命中率决定：$P(X=1)=p_{access}$，$P(X=0)=1-p_{access}$。

第四步：计算互信息。互信息$I(X;Y)$衡量观测$Y$对推断$X$的贡献。从直观理解：在观测之前，攻击者对$X$的不确定性是$H(X)$比特（$X$的熵）；观测到$Y$之后，剩余的不确定性是$H(X|Y)$比特（条件熵）。观测消除的不确定性即为获得的信息量：$I(X;Y) = H(X) - H(X|Y)$。将二值信道的概率参数代入，可推导出互信息的计算公式：

\begin{equation}
MI_{binary} = \sum_{x \in \{0,1\}} \sum_{y \in \{0,1\}} P(X=x, Y=y) \cdot \log_2 \frac{P(X=x, Y=y)}{P(X=x) \cdot P(Y=y)}
\end{equation}

其中联合概率$P(X=x, Y=y)$由$P(X=x) \cdot P(Y=y|X=x)$计算得到，边缘概率$P(Y=y)$由$\sum_x P(X=x, Y=y)$计算得到。

第五步：理解结果。当$f=0$（无噪声）时，观测结果$Y$完全确定输入$X$（命中对应访问，未命中对应未访问），此时互信息达到最大值$MI_{binary} = H(p_{access})$，即每次观测提供$H(p_{access})$比特信息。当$f>0$时，噪声使部分观测结果不可靠，互信息下降。当缓存刷新启用时，攻击者始终观测到未命中，观测结果与输入无关，互信息为零。

\paragraph{采样频率}

在本项目中，一次采样指的是攻击者完成一轮完整的Flush+Reload测量流程：用clflush刷新4张T表第0号缓存行→通知Victim加密一次→等待Victim完成→Reload测量4个缓存行的访问时间→记录命中/未命中状态和密文。因此$f_{sample}$就是单位时间内能完成上述完整轮次的次数：

\begin{equation}
f_{sample} = \frac{f_{CPU}}{C_{avg}}
\end{equation}

其中$f_{CPU}$为CPU主频（Hz），$C_{avg}$为单次采样的平均CPU周期数。$C_{avg}$通过实验测量：记录整个采样过程的起止时间戳（rdtsc），用总周期数除以样本数即得。以实验平台（i5-1135G7，2.4GHz）为例，AES-128基准配置下$C_{avg}$约17500周期，$f_{sample}$约137000次/秒。密钥越长，单次采样涉及的加密和等待开销越大，$C_{avg}$越高，$f_{sample}$越低，泄露带宽也随之下降。

\subsection{攻击成功率}

\begin{equation}
Success\_Rate = \frac{N_{success}}{N_{total}} \times 100\%
\end{equation}

其中$N_{total}$为总实验轮数，$N_{success}$为成功恢复完整密钥的轮数。

攻击成功率是最终的综合评估指标，受上述所有指标的间接影响。样本总数决定了排除法的收敛程度，未命中率决定了有效信息量，时间统计指标反映了测量质量，泄露带宽量化了信息泄露的速率。在实际评估中，攻击成功率与样本数呈单调递增关系：样本数不足时成功率接近0，超过阈值后迅速上升至接近100\%。该阈值受噪声水平和密钥长度的影响，噪声越高、密钥越长，所需样本数越多。泄露带宽则从另一个维度评估攻击威胁：即使攻击成功率相同，泄露带宽越高意味着攻击者在单位时间内能获取更多信息，攻击速度越快。因此，泄露带宽和攻击成功率从速率和可靠性两个互补的角度刻画了侧信道攻击的威胁程度。

\section{缓解措施实现}

本实验实现了三种针对Flush+Reload攻击的配置措施——CPU核心绑定、随机噪声注入和缓存刷新，可单独使用也可组合。

\subsection{CPU核心绑定}

实验中Attacker和Victim通过sched\_setaffinity绑定到同一CPU核心。同核运行时两个进程交替使用同一套L1d/L2缓存，缓存状态的变化直接反映受害者的访问行为，测量信号最清晰。由于x86架构的clflush指令会从所有缓存层级驱逐指定缓存行，核心绑定并非一种独立的缓解措施，本节仅将其作为实验的基本配置。

\subsection{随机噪声注入}

随机噪声注入的方法为：Victim加密完成后随机访问T表监控缓存行，干扰攻击者的观测。攻击者依赖"缓存行未被访问=未命中"的推理，如果Victim自身在加密以外额外访问了监控缓存行，攻击者的未命中观测即不再可靠。

噪声注入采用的是访问式注入（access-based）：Victim加密完成后，以概率$p_{noise}$对每张T表的第0号缓存行执行一次内存读取操作（mov指令），将该缓存行重新加载到CPU缓存中。低噪声$p_{noise} = 30\%$，中噪声$p_{noise} = 40\%$，高噪声$p_{noise} = 50\%$。读取操作使用普通的内存加载指令（mov），不使用clflush，读取目标为第0号缓存行内的随机表项（索引0-15）。

噪声注入对排除法的影响如下。当第0缓存行在加密中未被访问（真实未命中）时，噪声以概率$p_{noise}$将其重新加载到缓存中，使攻击者观测到命中（假阳性）。因此有效未命中率变为：

\begin{equation}
p'_{miss} = p_{miss} \times (1 - p_{noise})
\end{equation}

有效排除样本数从$N \cdot p_{miss}$减少为$N \cdot p_{miss} \times (1 - p_{noise})$。以AES-128为例，无噪声时$p_{miss} \approx 0.076$，有效排除样本约$5000 \times 0.076 = 380$；低噪声（$p_{noise} = 30\%$）时$p'_{miss} \approx 0.053$，有效排除样本约266；高噪声（$p_{noise} = 50\%$）时$p'_{miss} \approx 0.038$，有效排除样本约190。噪声越强，有效排除样本越少，收敛速度越慢，需要更多样本才能达到相同的密钥恢复置信度。

\subsection{缓存刷新}

缓存刷新是最有效的缓解措施：Victim加密完成后立即用clflush指令驱逐T表监控缓存行。实现方式为对4张$Te$表的第0号缓存行（即攻击者监控的目标缓存行）执行两次clflush指令，两次之间插入延迟和内存屏障（mfence、lfence），确保缓存行被彻底驱逐。

缓存刷新使攻击者的Reload阶段始终观测到缓存未命中，无论Victim是否访问过T表。此时命中率降为0，攻击者无法区分"Victim访问了该缓存行"和"Victim未访问该缓存行"两种情况，Flush+Reload攻击完全失效。

缓存刷新的性能开销：8次clflush（4个缓存行$\times$2次）加2次mfence、2次lfence，外加100次空循环延迟，按每次clflush约100周期、每次mfence和lfence约30周期估算，clflush与屏障开销约920周期，加上延迟循环，总开销约1000周期，约为一次AES-128加密（约200周期）的5倍。虽然开销远小于刷新全部128个缓存行的方案，但在对性能敏感的场景中仍不可忽视。缓存刷新适合作为安全测试的基准对照（验证攻击在最强防御下确实失效），而非实际部署的缓解方案。

\section{本章小结}

本章给出了AES Flush+Reload攻击测试框架的设计与实现。

攻击方法的核心是排除法。第一阶段通过监控T表第0缓存行的访问状态，利用密文与查表索引的代数关系逐步排除错误密钥候选，恢复最后一轮子密钥。本文进行了完整推导：单样本分析、多样本累积统计、期望样本量估计。仅监控第0缓存行的策略在信息量与测量效率之间实现了最优平衡。

第二阶段采用等效块解密方法，将同一批Flush+Reload监控数据复用于倒数第二轮密钥恢复。从密文出发，通过等效解密推导出解密中间值$IV_1$，然后对$Kd_1$的每个字节直接执行排除法——枚举S盒输出值，计算对应的$Kd_1$候选并排除。恢复$Kd_1$后，经MixColumns变换得到加密倒数第二轮子密钥。AES-128仅凭最后一轮密钥即可逆推原始密钥，无需第二阶段。AES-192和AES-256必须恢复倒数第二轮密钥才能完整逆推原始密钥。

系统实现层面，父子进程模型和管道同步通信协议保证了攻击流程的时序正确，量化评估指标体系涵盖基础指标、时间统计指标、泄露带宽和攻击成功率，综合衡量了攻击的数据需求、测量质量、信号强度和信息泄露速率，缓解措施实现了随机噪声注入和缓存刷新，其中缓存刷新能使攻击完全失效，但性能开销亦较大。
